\section{Experiments}


\subsection{Experiment Setup}
\noindent\textbf{Implementation details.}
We use CLIP VIT-L/14~\cite{radford2021learning} as the visual encoder. We uniformly sample 100 frames and encode each frame individually, representing each by the feature of the CLS token for efficiency. BEATs~\cite{chen2022beats} is employed as the audio encoder, and Whisper-Large-v2~\cite{radford2023robust} is used for speech encoding. Both audio and speech are processed from the waveforms of 5.12-second clips, with the number of tokens varying according to the video's duration. 
The audio and speech adapters are just randomly initialized. Furthermore, Vicuna-v1.5-7b~\cite{chiang2023vicuna} is adopted as the large language model. {\color{black} More details are in Appendix}.


\noindent\textbf{Evaluation metrics.}
Using our LongVALE test set, we evaluate models on three omni-modal fine-grained understanding tasks.
For omni-modal temporal video grounding, we report Recall@1 at IoU thresholds of \{0.3,0.5,0.7\} and mean IoU (mIoU).
For omni-modal dense video captioning, we assess caption quality using CIDEr~\cite{vedantam2015cider} and METEOR~\cite{banerjee2005meteor}, and employ SODA\_c~\cite{fujita2020soda} for overall story-level evaluation.
For omni-modal segment captioning, we use BLEU-4~\cite{papineni2002bleu}, ROGUE-L~\cite{lin2004automatic}, METEOR~\cite{banerjee2005meteor} and CIDEr~\cite{vedantam2015cider} for standard caption quality evaluation. 


\subsection{Main Results}
Table~\ref{tab:SOTA} presents the comparison results of our model with existing open-sourced video LLMs on the three omni-modal fine-grained tasks on our LongVALE test set, where we ensure optimal evaluation of the existing models.
Our LongVALE-LLM (7B) supports video, audio and speech input with fine-grained temporal understanding ability, and outperforms other video LLMs by a significant margin across all three tasks.
Despite VideoLLaMA~\cite{zhang2023video}, PandaGPT~\cite{su2023pandagpt} and NExT-GPT~\cite{wu2023next} also support audio-visual input, they are limited to processing a few video frames (\eg, 8 frames), resulting in poor performances on fine-grained, time-sensitive tasks. 
Besides, VTimeLLM~\cite{huang2024vtimellm} and TimeChat~\cite{ren2024timechat} can understand specific temporal events in videos, but they focus solely on visual events and fail to incorporate crucial audio information for a complete video understanding. 
Therefore, it is essential to integrate both audio and visual information with boundary-aware training on rich omni-modality data, \ie, LongVALE, to achieve precise video comprehension.

  
\begin{table}
  \centering
  \resizebox{0.9\linewidth}{!}{ 
  \begin{tabular}{l|c|cc}
    \toprule
    Method & \#Pairs & AVSD & Music-AVQA \\
    \midrule
    PandaGPT (13B)~\cite{su2023pandagpt} & 128M & 26.1 & 33.7 \\
    Macaw-LLM (7B)~\cite{lyu2023macaw} & 0.3M & 34.3 & 31.8 \\
    VideoLLaMA (7B)~\cite{zhang2023video} & 2.8M & 36.7 & 36.6 \\
    X-InstructBLIP (13B)~\cite{panagopoulou2023x} & 32M & - & 44.5 \\
    AV-LLM (13B)~\cite{shu2023audio} & 1.6M & 52.6 & 45.2 \\
    OneLLM (7B)~\cite{han2024onellm} & 1007M & - & 47.6 \\
    AVicuna (7B)~\cite{tang2024avicuna} & 1.1M & \underline{53.1} & \textbf{49.6} \\
    \midrule
    \rowcolor{gray!20} LongVALE-LLM (7B) (ours) & 0.7M & \textbf{54.8} & \underline{49.4}\\
    \bottomrule
  \end{tabular}}
  \vspace{-1.5mm}
  \caption{Comparison with existing LLM-based methods on open-ended audio-visual question answering benchmarks on a zero-shot setting. \# Pairs: the adopted instruction-response pairs.}
  \label{tab:AVQA}
  \vspace{-5mm}
\end{table}

\subsection{Zero-Shot Performance on General AVQA}
Besides the ability for fine-grained omni-modality video understanding tasks, we also explore whether our model trained on LongVALE can address a broader range of audio-visual questions in a zero-shot setting. 
We employ the AVSD~\cite{alamri2019audio} and Music-AVQA~\cite{li2022learning} benchmarks and conduct a GPT-assisted evaluation to assess the accuracy of the generated answers as same as the protocol~\cite{maaz2023video} used by other LLM-based methods shown in Tab.~\ref{tab:AVQA}.
We can observe that, despite using significantly less data, our model surprisingly achieves state-of-the-art performance on AVSD and highly competitive results on Music-AVQA.  
Notably, existing multi-modal LLMs all rely on large amounts of audio-related training data. 
For example, OneLLM~\cite{han2024onellm} and AVicuna~\cite{tang2024avicuna} use 460K and 350K audio/audio-visual instruction-response pairs for training, respectively. 
In contrast, we achieve even better results using only total 32.7K audio-visual samples from our LongVALE dataset, accounting for less than 10\% of the data they use. 
This clearly validates the robustness and generalization of our LongVALE-trained model, demonstrating that the comprehensive omni-modality captions in our dataset can effectively enhance the model's general cross-modal reasoning capability.


\begin{table}
  \centering
  \resizebox{\linewidth}{!}{ 
  \begin{tabular}{c|c|c|c|c}
    \toprule
    BP & IT &Omni-TVG$_{mIoU}$ & Omni-DVC$_{CIDEr}$ & Omni-SC$_{CIDEr}$  \\
     \midrule
     V~\cite{huang2024vtimellm} & \ding{55} &12.6 / 3.1 & 0.1 / 0.1 & 0.4 / 0.3  \\
     V~\cite{huang2024vtimellm} + Ours & \ding{55} & 25.6 / 7.0 & 7.3 / 7.0  & 21.1 / 17.5 \\
     % \midrule
     V~\cite{huang2024vtimellm} + Ours & V~\cite{huang2024vtimellm} & 11.7 / 3.8& 0.2 / 0.2 & 3.1 / 2.4  \\
     \rowcolor{gray!20} V~\cite{huang2024vtimellm} + Ours & V~\cite{huang2024vtimellm} + Ours & 26.0 / 7.3 &  7.8 / 7.7 & 25.1 / 19.4 \\
    \bottomrule
  \end{tabular}}
  \vspace{-1.5mm}
  \caption{Ablation study of the data used in different training stages. BP: boundary perception. IT: instruction tuning. V~\cite{huang2024vtimellm}: the visual-only data used in VTimeLLM~\cite{huang2024vtimellm}. 
  }
  \label{tab:ablation_stage}
  \vspace{-2mm}
\end{table}

\begin{table}
  \centering
  \resizebox{\linewidth}{!}{ 
  \begin{tabular}{c|c|c|c}
    \toprule
     Audio-Visual Correlation &Omni-TVG$_{mIoU}$ & Omni-DVC$_{CIDEr}$ & Omni-SC$_{CIDEr}$  \\
     \midrule
    \ding{55} & 23.7 / 6.5 & 3.5 / 3.4 & 10.4 / 10.2 \\
    \rowcolor{gray!20} \checkmark & 25.6 / 7.0 & 7.3 / 7.0  & 21.1 / 17.5 \\
    \bottomrule
  \end{tabular}}
  \vspace{-1.5mm}
  \caption{Ablation study of audio-visual correlation reasoning in captioning for the boundary perception training stage.}
  \label{tab:ablation_avc}
  \vspace{-2mm}
\end{table}

\begin{table}
  \centering
  \resizebox{0.9\linewidth}{!}{ 
  \begin{tabular}{c|c|c|c}
    \toprule
     Modality &Omni-TVG$_{mIoU}$ & Omni-DVC$_{CIDEr}$ & Omni-SC$_{CIDEr}$ \\
    \midrule
     V & 12.6 / 2.2 & 5.9 / 3.6 &14.6 / 11.7 \\
     V+A & 15.6 / 3.1 & 7.2 / 5.0 & 17.6 / 14.7 \\
     V+S & 15.2 / 2.9 & 6.7 / 5.2 & 17.3 / 14.6  \\
     \rowcolor{gray!20} V+A+S & 17.1 / 3.0 & 7.8 / 5.6 & 18.9 / 15.4 \\
    \bottomrule
  \end{tabular}}
  \vspace{-1.5mm}
  \caption{Ablation study of models trained on LongVALE with different modalities. V: vision. A: audio. S: speech. }
  \label{tab:ablation_modality}
  \vspace{-5mm}
\end{table}

\subsection{Ablation Study and Qualitative Results}
We provide detailed ablation studies about our training data, training stages, and modalities as follows.
Since LongVALE is a challenging benchmark, we split the test set into easy and hard subsets based on the ratio of average event duration to video length for more in-depth analysis. 
The easy subset ratios from 100\% to 9.3\% (585 videos), and the hard subset ratios from 9.3\% to 0.95\% (586 videos). 
Note that in Tab.~\ref{tab:ablation_stage}-\ref{tab:ablation_modality}, easy/hard subset results are on the left/right.

\noindent\textbf{Impact of LongVALE in different training stages.}
In Tab.~\ref{tab:ablation_stage}, we observe significant performance boosts across all tasks by adding our data in both boundary perception and instruction tuning stages. 
It demonstrates the essential role of our data in promoting omni-modality fine-grained video understanding.
Besides, the necessity of instruction tuning is evident when comparing the second and last rows. 
It shows consistent improvements across all three tasks, especially in segment captioning, indicating instruction tuning on high-quality dialog data further enhances comprehension and reasoning abilities for omni-modal events.

\noindent\textbf{Importance of audio-visual correlation (AVC).}
We simply concatenate the generated single-modality captions in Sec.~\ref{sec:3.3} as a naive caption for each omni-modal event boundary without AVC reasoning.
Then, we convert them into template-based dialogues for boundary perception tuning. 
The results in Tab.~\ref{tab:ablation_avc} show that the model trained with our omni-modality captions with AVC reasoning achieves significantly better performance, especially for captioning tasks, demonstrating the effectiveness of AVC to facilitate the model's capability of cross-modal reasoning.

\noindent\textbf{Impact of using different modalities.}
In Tab.~\ref{tab:ablation_modality}, we can see that adding audio or speech modality significantly improves the performances across all three tasks, with the best results achieved when all three modalities are used. 
This highlights the strength of omni-modality input for multi-modal video understanding.
Note that we use only our dataset during both training stages to provide a clearer comparison. 

\noindent\textbf{Qualitative results.}
In Fig.~\ref{fig:visualization}, VTimeLLM~\cite{huang2024vtimellm} with only visual data, misidentifies singing as raising hands. In contrast, our model integrates audio signals (\ie, man singing, crowd cheers) and performs cross-modal reasoning to accurately and comprehensively describe the event in the specific moment. Besides, our model effectively associates audio-visual information to provide correct answers for general AVQA. {\color{black}More examples can be found in Appendix.}



\begin{figure}[t]
  \centering
  \setlength{\abovecaptionskip}{1.0mm}
   \includegraphics[width=1.0\linewidth]{figures/new/fig5.jpg}
   \caption{Qualitative results. The orange text highlights audio-visual correlation for accurate and complete video understanding. Samples are from LongVALE and Music-AVQA test sets. 
   }
   \label{fig:visualization}
     \vspace{-5mm}
\end{figure}



\section{Conclusion}
We present a scalable pipeline to build LongVALE, the first benchmark for omni-modality fine-grained video understanding, featuring 105K omni-modal events with temporal boundaries and relation-aware captions. 
Benefiting from the dataset, our model exhibits distinct capabilities of both cross-modal reasoning and fine-grained temporal understanding that are absent in existing video LLMs, making a crucial step toward an intelligent video assistant.
In the future, we will expand our LongVALE with more high-quality data and advance the model's architecture to improve video semantic density and cross-modal interaction. 

